\documentclass[10pt,a4paper]{amsart}
\usepackage[margin=19mm]{geometry}
\usepackage{amsmath,amssymb,mathtools}
\usepackage[scr=rsfs,cal=euler]{mathalfa}
\usepackage{xfrac}
\usepackage[unicode,hidelinks,bookmarks=false]{hyperref}
\newcommand{\ie}{i.e.\ }
\newcommand{\cf}{cf.\ }
\newcommand{\resp}{respectively\ }
\newcommand{\Subsection}{\subsection}
\providecommand{\nobreakdash}{\nobreak\hbox{-}}
\setlength{\emergencystretch}{2em}
\multlinegap0pt
\usepackage{amssymb}
\newtheorem{lemma}{Lemma}
\newcommand{\R}{{\mathbb{R}}}
\newcommand{\T}{{\mathbb{T}}}
\newcommand{\Z}{{\mathbb{Z}}}
\title{A priori bounds and equicontinuity of orbits\\for the intermediate long wave equation}
\author{Benjamin Harrop-Griffiths, Rowan Killip, Monica Visan}
\date{Source: arXiv:2506.23868v2 (CC BY 4.0)}
\begin{document}
\maketitle

\noindent\emph{Source and licence.} A priori bounds and equicontinuity of orbits\\for the intermediate long wave equation. Version arXiv:2506.23868v2 (CC BY 4.0), distributed by arXiv under the Creative Commons Attribution 4.0 International licence, http://creativecommons.org/licenses/by/4.0/, as recorded in the arXiv licence metadata for this exact version and shown on the abstract page https://arxiv.org/abs/2506.23868v2. Full licence notice: \texttt{licenses/arxiv-2506.23868-CC-BY-4.0.txt}.\par\smallskip

\noindent\textit{Editorial scope.} Counted unit: the article's lemma L:HS (source0.tex lines 265--302). The article's complete original proof of it is delivered with it (source0.tex lines 303--387, one \texttt{proof} environment of 85 lines). No mathematics is written, repaired or re-derived: the statement and the proof are the article's own lines, in the article's own notation, and no retained line is substituted or re-flowed.  Retained uncounted as the source-local context the proof needs: lines 168--170.  Nothing was omitted from any retained span, so the package carries no omission record.  No retained line contains a citation, so the wrapper carries no bibliography and no bibliography entry.  No retained line contains a figure environment, an image-inclusion command or drawing code, so no image is delivered and the package declares no asset dependency.  Editorial formatting only, in addition: an AMS \texttt{amsart} preamble that restates the article's own declarations and macros used by the retained lines, namely the environments \texttt{lemma} on separate counters, together with the standard packages those lines need.  The retained environments print in this document as Lemma 1 in order of appearance, whereas the article prints the same items with the numbering of its own full document.  The complete native source of the article as distributed in the arXiv e-print for this exact version is bundled unchanged as \texttt{sources/180743/source0.tex}.
\begin{equation}\label{symbol}
a_h(\xi) = \xi+ \tfrac1{2h}(e^{-2h\xi}-1).
\end{equation}

\begin{lemma}\label{L:HS}
For $h,\kappa>0$, we have
\begin{align}\label{I2}
 \bigl\| \sqrt{R_0(\kappa)}\, q\, \sqrt{R_0(\kappa)} \bigr\|_{\mathfrak I_2(\mathcal M)}^2
 = \begin{cases}\displaystyle\int_{\R}  F(\xi;\kappa,h)\, |\widehat q(\xi)|^2\,d\xi &\quad\text{if \(\mathcal M = \R\),}\vspace{1em}\\ \displaystyle\sum_{\xi\in2\pi\Z} F(\xi;\kappa,h)\, |\widehat q(\xi)|^2&\quad\text{if \(\mathcal M = \T\),}\end{cases}
\end{align}
where 
\begin{align}\label{F defn}
 F(\xi;\kappa,h) &:= \begin{cases}
 	\displaystyle\ \int_{\R} \bigl[a_h(\eta) +\kappa\bigr]^{-1}\bigl[a_h(\xi+\eta) + \kappa\bigr]^{-1} \,\tfrac{d\eta}{2\pi}  & \quad\text{on $\R$,}\\[2ex]
	\displaystyle\sum_{\eta\in2\pi\Z} \bigl[a_h(\eta) +\kappa\bigr]^{-1}\bigl[a_h(\xi+\eta) + \kappa\bigr]^{-1} & \quad\text{on $\T$,}
\end{cases}
\end{align}
with \(a_h(\xi)\) defined as in \eqref{symbol}.

The function \(F(\xi;\kappa,h)\) is smooth, non-negative, even in \(\xi\), and satisfies the estimate
\begin{equation}\label{F main bound}
F(\xi;\kappa,h)\simeq \bigl[\tfrac{h\xi^2}{1+h|\xi|} +\kappa\bigr]^{-1}\Bigl[\sqrt{\tfrac{1+h\kappa}{h\kappa}} + \log\bigl(1 + \tfrac{h|\xi|}{1+h\kappa}\bigr)\Bigr].
\end{equation} 

In particular, if $-\frac12<s\leq 0$ then
\begin{align}\label{HS bdd}
 \| \sqrt{R_0(\kappa)}\, q\, \sqrt{R_0(\kappa)}\|_{\mathfrak I_2(\mathcal M)}^2 \lesssim_s\begin{cases} \kappa^{-(1-2|s|)} \|q\|_{H^s_\kappa(\mathcal M)}^2&\quad\text{if \(\kappa>\tfrac1h\),}\vspace{1em}\\\bigl(\tfrac \kappa h\bigr)^{-(\frac32-|s|)}\bigl\|\tfrac qh\bigr\|_{H_{\sqrt{\kappa/h}}^s(\mathcal M)}^2&\quad\text{if \(0<\kappa\leq\tfrac1h\).}\end{cases}
\end{align}

Consequently, there exists a constant \(C_s>0\) so that if \(h,\delta>0\) and
\begin{equation}\label{k bigger}
\kappa \geq \max\Bigl\{\tfrac1h,\bigl[ 1 + \tfrac1{\delta^2} C_s \| q\|_{H^s}^2 \bigr]^{\frac{1}{1-2|s|}}\Bigr\}
\end{equation}
or
\begin{equation}\label{shallow k bigger}
\bigl[ 1 + \tfrac1{\delta^2} C_s\bigl\|\tfrac qh\bigr\|_{H^s}^2 \bigr]^{\frac{1}{\frac32-|s|}}\leq \tfrac\kappa h\leq \tfrac1{h^2}
\end{equation}
then
\begin{align}\label{small norm}
 \| \sqrt{R_0(\kappa)}\, q\, \sqrt{R_0(\kappa)}\|_{\mathfrak I_2(\mathcal M)}^2 <\delta^2.
\end{align}
\end{lemma}

\begin{proof}
The identities \eqref{I2} and \eqref{F defn} follow from considering the Fourier kernel of the operator \(\sqrt{R_0(\kappa)}\, q\, \sqrt{R_0(\kappa)}\).

Making the change of variables \(\eta\mapsto \eta-\xi\) in \eqref{F defn} proves that the map \(\xi\mapsto F(\xi;\kappa,h)\) is even.  That $F$ is also smooth and non-negative follows from the fact that the function $a_h$ is smooth and non-negative.

The estimate \eqref{HS bdd} follows readily from \eqref{F main bound}, while \eqref{small norm} follows from \eqref{HS bdd} under the assumptions \eqref{k bigger} and \eqref{shallow k bigger}.  Thus, it remains to estimate the size of $F(\xi;\kappa,h)$.  We present the details on the real line; the periodic setting can be treated analogously, with integrals being replaced by sums.  Further, as \(F\) is even in \(\xi\) and $|\widehat q(\xi)| = |\widehat q(-\xi)|$ (since $q$ is real-valued), it suffices to consider the case that \(\xi\geq0\). 

Denoting \(a(\xi) = a_1(\xi) = \xi + \tfrac12(e^{-2\xi}-1)\), we compute that \(a'(\xi) = 1 - e^{-2\xi}\) and \(a''(\xi) = 2e^{-2\xi}>0\), so \(a(\xi)\) is convex with a global minimum at \(\xi = 0\). By considering the asymptotic behavior of \(a(\xi)\) as \(\xi\to0\) and \(\xi\to \pm\infty\), we see that
\begin{equation}\label{a bounds}
a(\xi)\gtrsim \tfrac{\xi^2}{1+|\xi|}\quad\text{for all \(\xi\in \R\)}\quad\text{and}\quad a(\xi)\simeq \tfrac{\xi^2}{1+\xi}\quad\text{for all \(\xi\geq0\).}
\end{equation}

As \(a_h(\xi) = \frac1ha(h\xi)\), the estimate \eqref{a bounds} yields the (uniform in \(h\)) upper bound
\begin{align}
F(\xi;\kappa,h) &\lesssim \int_\R \bigl[\tfrac{h\eta^2}{1+h|\eta|} +\kappa\bigr]^{-1}\bigl[\tfrac{h(\xi+\eta)^2}{1+h|\xi+\eta|} + \kappa\bigr]^{-1} \,d\eta\notag\\
&\lesssim \int_{-\frac12\xi}^{\frac12\xi} \bigl[\tfrac{h\eta^2}{1+h|\eta|} +\kappa\bigr]^{-1}\bigl[\tfrac{h(\xi+\eta)^2}{1+h|\xi+\eta|} + \kappa\bigr]^{-1} \,d\eta\label{F upper bound}\\
&\quad  + \int_{\frac12\xi}^\infty \bigl[\tfrac{h\eta^2}{1+h|\eta|} +\kappa\bigr]^{-1}\bigl[\tfrac{h(\xi+\eta)^2}{1+h|\xi+\eta|} + \kappa\bigr]^{-1} \,d\eta,\notag
\end{align}
where the second inequality follows from the fact that the map
\[
\eta\mapsto \bigl[\tfrac{h\eta^2}{1+h|\eta|} +\kappa\bigr]^{-1}\bigl[\tfrac{h(\xi+\eta)^2}{1+h|\xi+\eta|} + \kappa\bigr]^{-1}
\]
is symmetric about \(\eta = -\frac12\xi\).


For the first summand on RHS\eqref{F upper bound}, we estimate
\begin{align*}
\int_{-\frac12\xi}^{\frac12\xi}\bigl[\tfrac{h\eta^2}{1+h|\eta|} +\kappa\bigr]^{-1}\bigl[\tfrac{h(\xi+\eta)^2}{1+h|\xi+\eta|} + \kappa\bigr]^{-1} \,d\eta \simeq \bigl[\tfrac{h\xi^2}{1+h|\xi|} +\kappa\bigr]^{-1}\int_0^{\frac12\xi}\bigl[\tfrac{h\eta^2}{1+h|\eta|} +\kappa\bigr]^{-1}\,d\eta.
\end{align*}
We then consider several cases depending on the relative size of \(\xi,\kappa,h\).

First, suppose that \(\kappa>\frac1h\). If additionally \(0\leq\xi\leq \kappa\), then
\begin{align*}
\int_0^{\frac12\xi}\bigl[\tfrac{h\eta^2}{1+h|\eta|} +\kappa\bigr]^{-1}\,d\eta&\simeq \int_0^{\frac12\xi}\tfrac1\kappa\,d\eta\simeq \tfrac\xi\kappa\lesssim 1,
\end{align*}
whereas if \(\xi>\kappa\) we have
\begin{align*}
\int_0^{\frac12\xi}\bigl[\tfrac{h\eta^2}{1+h|\eta|} +\kappa\bigr]^{-1}\,d\eta\simeq \int_0^{\frac12\kappa}\tfrac1\kappa\,d\eta + \int_{\frac12\kappa}^{\frac12\xi}\tfrac1\eta\,d\eta\simeq 1 + \log\bigl(\tfrac\xi\kappa\bigr).
\end{align*}

Second, suppose that \(0<\kappa\leq \frac1h\). If \(0\leq\xi\leq \sqrt{\frac\kappa h}\) then
\[
\int_0^{\frac12\xi}\bigl[\tfrac{h\eta^2}{1+h|\eta|} +\kappa\bigr]^{-1}\,d\eta \simeq \int_0^{\frac12\xi}\tfrac1\kappa\,d\eta \simeq \tfrac\xi\kappa\lesssim \tfrac1{\sqrt{h\kappa}},
\]
if \(\sqrt{\frac\kappa h}<\xi\leq\frac1h\) then
\begin{align*}
\int_0^{\frac12\xi}\bigl[\tfrac{h\eta^2}{1+h|\eta|} +\kappa\bigr]^{-1}\,d\eta&\simeq \int_0^{\frac12\sqrt{\frac\kappa h}}\tfrac1\kappa\,d\eta + \int_{\frac12\sqrt{\frac\kappa h}}^{\frac12\xi}\tfrac1{h\eta^2}\,d\eta\simeq \tfrac1{\sqrt{h\kappa}},
\end{align*}
and if \(\xi>\frac1h\) then
\begin{align*}
\int_0^{\frac12\xi}\bigl[\tfrac{h\eta^2}{1+h|\eta|} +\kappa\bigr]^{-1}\,d\eta&\simeq \int_0^{\frac12\sqrt{\frac\kappa h}}\tfrac1\kappa\,d\eta + \int_{\frac12\sqrt{\frac\kappa h}}^{\frac1{2h}}\tfrac1{h\eta^2}\,d\eta + \int_{\frac1{2h}}^{\frac12\xi}\tfrac1\eta\,d\eta\simeq \tfrac1{\sqrt{h\kappa}} + \log(h\xi).
\end{align*}

Putting these cases together, we arrive at the estimate
\[
\int_0^{\frac12\xi} \bigl[\tfrac{h\eta^2}{1+h|\eta|} +\kappa\bigr]^{-1} \,d\eta\lesssim \sqrt{\tfrac{1+h\kappa}{h\kappa}} + \log\bigl(1 + \tfrac{h\xi}{1+h\kappa}\bigr),
\]
which yields an acceptable upper bound for the first summand on RHS\eqref{F upper bound}. We also note that these estimates prove that
\begin{equation}\label{lower bound helper}
\int_0^{\frac12\max\{\xi,\kappa,\sqrt{\frac\kappa h}\}} \bigl[\tfrac{h\eta^2}{1+h|\eta|} +\kappa\bigr]^{-1} \,d\eta\simeq \sqrt{\tfrac{1+h\kappa }{h\kappa}} + \log\bigl(1 + \tfrac{h\xi}{1+h\kappa}\bigr),
\end{equation}
which we will use in our proof of the lower bound below.

Turning to the second summand on RHS\eqref{F upper bound}, we estimate
\[
\int_{\frac12\xi}^\infty \bigl[\tfrac{h\eta^2}{1+h|\eta|} +\kappa\bigr]^{-1}\bigl[\tfrac{h(\xi+\eta)^2}{1+h|\xi+\eta|} + \kappa\bigr]^{-1} \,d\eta\simeq \int_{\frac12\xi}^\infty \bigl[\tfrac{h\eta^2}{1+h|\eta|} +\kappa\bigr]^{-2}\,d\eta.
\]
We proceed as for the first summand: When \(\kappa>\frac1h\), we consider the cases \(\xi>\kappa\) and \(0\leq\xi\leq \kappa\) to obtain the estimate
\[
\int_{\frac12\xi}^\infty \bigl[\tfrac{h\eta^2}{1+h|\eta|} +\kappa\bigr]^{-2}\,d\eta\lesssim \bigl[\tfrac{h\xi^2}{1+h\xi} +\kappa\bigr]^{-1}\Bigl[\sqrt{\tfrac{1+h\kappa}{h\kappa}} + \log\bigl(1 + \tfrac{h\xi}{1+h\kappa}\bigr)\Bigr].
\]
Similarly, when \(0<\kappa\leq \frac1h\) we consider the cases \(\xi>\frac1h\), \(\sqrt{\frac\kappa h}<\xi\leq \frac1h\), and \(0\leq\xi\leq \sqrt{\frac\kappa h}\) to bound
\[
\int_{\frac12\xi}^\infty \bigl[\tfrac{h\eta^2}{1+h|\eta|} +\kappa\bigr]^{-2}\,d\eta\lesssim \bigl[\tfrac{h\xi^2}{1+h\xi} +\kappa\bigr]^{-1}\Bigl[\sqrt{\tfrac{1+h\kappa }{h\kappa}} + \log\bigl(1 + \tfrac{h\xi}{1+h\kappa}\bigr)\Bigr].
\]

Combining our estimates for the two summands on RHS\eqref{F upper bound} and recalling that \(F\) is even in \(\xi\) gives us the upper bound in \eqref{F main bound}.

Turning to the lower bound in \eqref{F main bound}, we take \(\xi\geq 0\) and combine \eqref{F defn} and \eqref{a bounds} to obtain
\begin{align*}
F(\xi;\kappa,h) &\gtrsim \int_0^{\frac12\max\{\xi,\kappa,\sqrt{\frac\kappa h}\}} \bigl[\tfrac{h\eta^2}{1+h|\eta|} +\kappa\bigr]^{-1}\bigl[\tfrac{h(\xi+\eta)^2}{1+h(\xi+\eta)} + \kappa\bigr]^{-1} \,d\eta\\
&\gtrsim \bigl[\tfrac{h\xi^2}{1+h\xi} + \kappa\bigr]^{-1}\int_0^{\frac12\max\{\xi,\kappa,\sqrt{\frac\kappa h}\}} \bigl[\tfrac{h\eta^2}{1+h|\eta|} +\kappa\bigr]^{-1} \,d\eta.
\end{align*}
Applying \eqref{lower bound helper} and again using that \(F\) is even in \(\xi\) completes the proof of \eqref{F main bound}. 
\end{proof}

\end{document}
